\begin{abstract}
\noindent
Single-image 3D scene generation has advanced rapidly, but most systems stop at a
viewable, inert result: a panorama, a layered mesh, or a radiance field that
supports novel-view synthesis but has no motion and nothing to interact with.
We present \textbf{Into Frame}, a system that turns one photograph into an
\emph{inhabitable} environment, in which discrete objects exist as individual
simulated assets, ground cover and vegetation respond to contact and wind, and
the scene is lit by a physically based estimate of the original illumination, all
at real-time VR frame rates.
Our central design commitment is to reconstruct each part of the scene by a
method suited to its physical role, rather than extracting everything from one
representation: terrain, objects, ground cover, water and sky are handled
separately and by different means.
We outline the pipeline built on this principle and focus on four contributions:
a projection from calibrated panoramic depth to a top-down height field that
preserves each cell's observed panorama pixel; a terrain reconstruction that
keeps visual similarity with the original image while applying fluvial erosion to
give plausible relief where nothing was observed; a pattern-texturing scheme that
avoids the projective smearing an equirectangular photograph produces on the
ground; and a point-process population synthesis that scatters vegetation and
objects to match the spatial statistics of the real detections.
We evaluate the system on several real-world images and report both what it does
well and a set of failure modes we were able to measure precisely. We argue that
the most transferable result is methodological: at this scale, most defects are
not model failures but \emph{assumption} failures, where a step valid for one
scene category is silently applied to another.
\end{abstract}
